\scanpage{50}
% archive sheet number: 49
\mspage{2 - 8 -}\hfill n\textsuperscript{o} 392

de $\underline{A}$ au moyen de $\handwritten{\Delta}$, le groupe $\mathrm{SL}(n,K)$ est un sous-groupe de $\mathrm{SL}(n,\underline{A})$.

PROPOSITION 4.- \uline{Le groupe} $\mathrm{SL}(n,K)$ \uline{est dense dans} $\mathrm{SL}(n,\underline{A})$.

Soit $G$ l'adhérence de $\mathrm{SL}(n,K)$ dans $\mathrm{SL}(n,\underline{A})$ ; comme $K$ est dense dans $\underline{A}$, (prop. 3), $G$ contient toutes les matrices de la forme $\underline{I}+\underline{a}.\underline{E}_{ij}$ pour $i\handwritten{\ne}j$ et $\underline{a}\handwritten{\in}\underline{A}$. Pour tout $\unsure{\underset{=}{p}|\underline{p}|p}\handwritten{\in}P$ et tout $\handwritten{\lambda\in}\overset{\handwritten{\scriptscriptstyle\wedge}}{K}_{\underset{=}{p}}$, soit $\handwritten{\lambda}(\underset{=}{p})$ l'adèle restreint $\underline{x}=(x_{\underset{=}{q}})_{\underset{=}{q}\handwritten{\in}P}$ tel que $x_{\underset{=}{p}}=\handwritten{\lambda}$ et $x_{\underset{=}{q}}=0$ pour $\underset{=}{q}\handwritten{\ne}\underset{=}{p}$ ; ce qui précède montre que $G$ contient les matrices $\underline{I}+\handwritten{\lambda}(\underset{=}{p})\underline{E}_{ij}$ pour $i\handwritten{\ne}j$. Mais, on sait que les matrices de la forme $\underline{I}+\handwritten{\lambda}\underline{E}_{ij}$, pour $\handwritten{\lambda\in}\overset{\handwritten{\scriptscriptstyle\wedge}}{K}_{\underset{=}{p}}$, engendrent le groupe $\mathrm{SL}(n,\overset{\handwritten{\scriptscriptstyle\wedge}}{K}_{\underset{=}{p}})$, (\uline{Alg}., chap. III, 3e éd., § \hspace{1em}, n\textsuperscript{o} \hspace{1em}, \hspace{1.7em}). Pour toute matrice $\underline{U}\handwritten{\in}\mathrm{SL}(n,\underline{A})$ désignons par $\underline{U}_{\underset{=}{p}}$ l'image canonique de $\underline{U}$ dans $\mathrm{SL}(n,\overset{\handwritten{\scriptscriptstyle\wedge}}{K}_{\underset{=}{p}})$ ; on voit donc que pour tout $\underset{=}{p}\handwritten{\in}P$, $G$ contient les matrices $\underline{U}\handwritten{\in}\mathrm{SL}(n,\underline{A})$ telles que $\underline{U}_{\underset{=}{q}}=\underline{I}$ pour tout $\underset{=}{q}\handwritten{\ne}\underset{=}{p}$. Puisque $G$ est un groupe, il contient aussi toutes les matrices $\underline{U}\handwritten{\in}\mathrm{SL}(n,\underline{A})$ telles que $\underline{U}_{\underset{=}{p}}=\underline{I}$ sauf pour un nombre \uline{fini} de $\underset{=}{p}\handwritten{\in}P$ ; or, la définition de la topologie de $\underline{A}$ montre aussitôt que l'ensemble de ces matrices est dense dans $\mathrm{SL}(n,\underline{A})$.

5. \uline{Le théorème de Krull-Akizuki}.

LEMME 1.- \uline{Soient} $A$ \uline{un anneau noethérien intègre dans lequel tout idéal premier non nul est maximal}, \uline{et} $M$ \uline{un} $A$-\uline{module de torsion de type fini}. \uline{Alors la longueur} $\handwritten{\ell}_A(M)$ \uline{de} $M$ \uline{est finie}.

En effet, comme $M$ est un module de torsion, tout idéal premier associé à $M$ est $\handwritten{\ne}(0)$, donc maximal. Le lemme résulte alors du chap. IV, § 2, n\textsuperscript{o} 5, prop. 7.

LEMME 2.- \uline{Soient} $A$ \uline{un anneau}, $T$ \uline{un} $A$-\uline{module}, $(T_{\handwritten{\unsure{\iota|\tau}}})$ \uline{une famille filtrante croissante de sous-modules de} $T$, \uline{de réunion} $T$. \uline{Alors} $\handwritten{\ell}_A(T)=\sup_{\handwritten{\iota}}(\,\handwritten{\ell}_A(T_{\handwritten{\iota}}))$.

On a $\handwritten{\ell}_A(T_{\handwritten{\iota}})\handwritten{\leq\ell}_A(T)$ pour tout $\handwritten{\iota}$. Le lemme est évident si aucun entier ne majore les $\handwritten{\ell}_A(T_{\handwritten{\iota}})$, les deux membres étant alors infinis. Sinon, soit $\handwritten{\iota}_0$ un indice pour lequel $\handwritten{\ell}_A(T_{\handwritten{\iota}})$ prenne sa plus grande valeur ; on a $T_{\handwritten{\iota}_0}=T$ puisque la famille
